\scanpage{46}
EGA IV 4\textordmasculine\quad Errata et Compléments.

\deleted{$\unsure{0|O}_{\mathrm{III}}$ 10.2.1. p.18 l. -6, au lieu de \textquotedbl en outre si \textquotedbl\ lire \textquotedbl si en outre\textquotedbl}

\deleted{III 4.4.12. , p.138 l.10 , au lieu de Chap V lire Chap IV 8.11.2.}

\deleted{III 3.4. p.119, supprimer les lignes 2 et 3, qui contiennent une promesse
qui ne sera pas tenue.}

% Bracket and surrounding line by hand.
\margin{\handwritten{?}}
III 4.3.6. , p. 133 l. 7 , supprimer \deleted{la relation} \textquotedbl; on a donc $K'=K\unsure{\otimes|\mathrm O}_{A}A'$ \textquotedbl,
qui est manifestement faux.

\deleted{$0_{\mathrm{IV}}$ 22.1.1. p.183 l. 12, remplacer (20.6.20) par (20.6.4) .}

\deleted{IV Sommaire p.223, l. 16 au lieu de \textquotedbl de présentation finie\textquotedbl\ lire
\textquotedbl localement de présentation finie\textquotedbl\ ; ligne 19, après \textquotedbl comme\textquotedbl\ ajouter
\textquotedbl étant\textquotedbl.}

\deleted{IV 3.1.10. \deleted{Ajouter à la proposition une partie (ii), où on} Supposer
\inserted{d'abord} seulement $f$ quasi-compact et quasi-séparé, et \deleted{où on} conclu\deleted{\unsure{t}}re que
$\mathrm{Ass}\,f_{\handwritten{*}}(\underline F)\handwritten{\subset} f(X)$, et même $\mathrm{Ass}\, f_{\handwritten{*}}(\underline F)\handwritten{\subset} f_{\handwritten{\unsure{*|\leftarrow}}}(\mathrm{Ass}\,\underline F)$ lorsque $X$ est localement noethérien; enfin, si de plus $f$ est entier, on a l'égalité dans
cette dernière inclusion.}

\deleted{IV 3 , à la page 43, ajouter après 3.2.8. un\deleted{\unsure{e}}}

\deleted{\uline{Corollaire} 3.2.9. Soit $\underline F$ un Module cohérent sur le préschéma localement
noethérien $X$. Pour toute partie \deleted{discrète} \inserted{discrète} $S$ de $X$, \inserted{i.e. loc. finie et} telle que deux points
distincts de $S$ ne so\deleted{\unsure{nt}}ient pas spécialisation l'un de l'au\deleted{\unsure{d}}tre, et \inserted{pour} tout sys\unsure{t}tème $(\underline G(s))_{s\;S}$ de Modules quotients de lo\deleted{\unsure{gg}}ngueur finie de $\underline F_s$ , considérons le Module quotient $\underline G=\underline F/\underline H$ de $\underline F$, où $\underline H$ est le noyau de l'homomorphisme \deleted{naturel} composé $\underline F\quad i_s\ (\underline F_s)\quad i_s\ (G(s))$ , $i_s\!:\operatorname{Spec} k(s)\quad S$
étant le monomorphisme canonique. Alors $\underline G$ est un Module quotient cohérent
de $\underline F$ sans cycles associés immergés. L'application qui, à tout système
$(S,(G(s))_{s\;S})$ comme ci-dessus, associé $\underline G$ , est une bijection de l'ensemble de ces systèmes sur l'ensemble des Modules quotients cohérents de}
% archive sheet number: 45
